\documentclass[11pt,a4paper]{article}
% =====================================================================
%  Gemeinsame Vorlage der Arbeitsblätter "Rekursive Datenstrukturen"
%  – Informatik 12 (gleiches Layout wie informatik/13/formale-sprachen/
%  arbeitsblaetter/ab-vorlage.tex, dazu Listings und Objektdiagramme)
%  Einbinden in jedem Blatt direkt nach \documentclass: % =====================================================================
%  Gemeinsame Vorlage der Arbeitsblätter "Rekursive Datenstrukturen"
%  – Informatik 12 (gleiches Layout wie informatik/13/formale-sprachen/
%  arbeitsblaetter/ab-vorlage.tex, dazu Listings und Objektdiagramme)
%  Einbinden in jedem Blatt direkt nach \documentclass: % =====================================================================
%  Gemeinsame Vorlage der Arbeitsblätter "Rekursive Datenstrukturen"
%  – Informatik 12 (gleiches Layout wie informatik/13/formale-sprachen/
%  arbeitsblaetter/ab-vorlage.tex, dazu Listings und Objektdiagramme)
%  Einbinden in jedem Blatt direkt nach \documentclass: \input{ab-vorlage}
%  Übersetzen mit XeLaTeX, LuaLaTeX, pdfLaTeX oder Tectonic.
%
%  Blatt:   xelatex ab-1-2-interfaces.tex
%  Lösung:  xelatex ab-1-2-interfaces-lsg.tex
%           (Datei enthält nur: \def\mitloesung{}\input{ab-1-2-interfaces};
%            füllt alle Lücken rot aus; siehe \lsg, \lsglinien, \lsgoder,
%            \lsgtext, lsgbild, \kreuz)
% =====================================================================
\usepackage[a4paper,left=16mm,right=16mm,top=23mm,bottom=13mm,headheight=12mm,headsep=3mm,footskip=6mm]{geometry}
\usepackage{iftex}
\ifPDFTeX
  \usepackage[T1]{fontenc}
  \usepackage[utf8]{inputenc}
  \usepackage{helvet}
  \usepackage[scaled=0.86]{beramono}
  \renewcommand{\familydefault}{\sfdefault}
\else
  \usepackage{fontspec}
  \setmainfont{texgyreheros}[Extension=.otf,UprightFont=*-regular,BoldFont=*-bold,ItalicFont=*-italic,BoldItalicFont=*-bolditalic]
  \setmonofont{DejaVuSansMono}[Extension=.ttf,UprightFont=*,BoldFont=*-Bold,ItalicFont=*-Oblique,BoldItalicFont=*-BoldOblique,Scale=0.86]
\fi
\usepackage[ngerman,shorthands=off]{babel}
\usepackage{amssymb}
\usepackage{xcolor,graphicx,tabularx,array,enumitem,fancyhdr,tikz,listings,multicol,calc,etoolbox,pifont}
\usepackage[most]{tcolorbox}
\usetikzlibrary{arrows.meta,positioning,calc,decorations.pathreplacing,shapes.multipart}
\usepackage[hidelinks]{hyperref}
\setlength{\parindent}{0pt}
\setlength{\parskip}{3pt}
\setlist{nosep,leftmargin=*}

% ---------- Lösungsschalter ----------
\newif\ifloesung
\ifdefined\mitloesung\loesungtrue\fi
\newcommand{\lsgfarbe}{\color{red!75!black}}

% ---------- Kopf- und Fußzeile (einheitliche Form) ----------
\newcommand{\lehrkraft}{Hau}
\newcommand{\themenbereich}{1. Rekursive Datenstrukturen}
\newcommand{\abnummer}{}
\pagestyle{fancy}\fancyhf{}
\renewcommand{\headrulewidth}{0.4pt}
\fancyhead[L]{\small Klasse 12\\Datum:}
\fancyhead[C]{\small Themenbereich:\\\themenbereich}
\fancyhead[R]{\small \lehrkraft\ifloesung\\{\lsgfarbe\bfseries Lösung}\fi}
\fancyfoot[C]{\scriptsize Informatik 12\,·\,Blatt \abnummer\ifloesung\ (Lösung)\fi\,·\,Seite \thepage}
\newcommand{\abtitel}[2]{\renewcommand{\abnummer}{#1}\begin{center}\Large\bfseries #1\quad\underline{#2}\end{center}\vspace{-2mm}}

% ---------- Boxen ----------
\newenvironment{aufgabe}[2][]{\begin{tcolorbox}[enhanced,breakable,colback=white,colframe=black,boxrule=0.7pt,arc=3mm,left=2mm,right=2mm,top=1mm,bottom=1.2mm,before skip=2mm,after skip=2mm]\textbf{Aufgabe #2)}\ifstrempty{#1}{}{ \textit{#1}}\par\vspace{0.3mm}}{\end{tcolorbox}}
\newtcolorbox{merke}[1][Merke]{enhanced,breakable,colback=green!5,colframe=green!40!black,boxrule=0.8pt,arc=2mm,left=2mm,right=2mm,top=1mm,bottom=1.2mm,before skip=2mm,after skip=2mm,title={#1},fonttitle=\bfseries,coltitle=white,toptitle=0.3mm,bottomtitle=0.3mm}
\newtcolorbox{hinweis}{enhanced,breakable,colback=yellow!12,colframe=orange!70!black,boxrule=0.6pt,arc=2mm,left=2mm,right=2mm,top=1mm,bottom=1mm,before skip=2mm,after skip=2mm}
\newcommand{\web}[1]{{\small\color{gray}\textit{Website: #1}}}

% ---------- Lücken, Linien, Karos ----------
\newcommand{\luecke}[1][3.5cm]{\rule[-0.5ex]{#1}{0.4pt}}
\newcommand{\linien}[1]{\par\vspace{1.5mm}\foreach \i in {1,...,#1}{\noindent\rule[0pt]{\linewidth}{0.3pt}\par\vspace{5.5mm}}\vspace{-3mm}}
\newcommand{\kariert}[1]{\par\vspace{1.5mm}\noindent\begin{tikzpicture}\draw[step=5mm,gray!55,thin] (0,0) grid (\linewidth-1pt,#1*5mm);\end{tikzpicture}\par}
\newcommand{\karobox}[2]{\begin{tikzpicture}\draw[step=5mm,gray!55,thin] (0,0) grid (#1,#2);\end{tikzpicture}}
\newcommand{\code}[1]{\texttt{#1}}
% Schritt-Nummern wie im Code und auf der Website: \nr{1} = ①
\newcommand{\nr}[1]{\ding{\numexpr171+#1\relax}}

% Lücke mit Lösung: \lsg[Mindestbreite]{Lösungstext}
%  Die Lösung steht (ohne eigene Höhe/Tiefe) auf derselben Linie wie im Blatt: Das Layout verschiebt sich nicht,
%  solange der Text nicht breiter als die Mindestbreite ist.
\newlength{\lsgbreite}
\newcommand{\lsg}[2][3.5cm]{\leavevmode\ifloesung\settowidth{\lsgbreite}{\,#2\,}\ifdim\lsgbreite<#1\setlength{\lsgbreite}{#1}\fi\rlap{\smash{\raisebox{0.8pt}{\makebox[\lsgbreite][l]{\lsgfarbe\,#2}}}}\luecke[\lsgbreite]\else\luecke[#1]\fi}
% Schreiblinien mit Lösung: \lsglinien{Anzahl}{Lösungstext} (gleiche Höhe wie \linien:
%  je Linie \baselineskip + 5,5 mm + \parskip, ausgemessen)
\newlength{\lsglinienhoehe}
\newcommand{\lsglinien}[2]{\ifloesung\par\vspace{1.5mm}\setlength{\lsglinienhoehe}{\dimexpr #1\dimexpr\baselineskip+5.5mm+\parskip\relax-16.25pt\relax}\noindent\begin{minipage}[t][\lsglinienhoehe][t]{\linewidth}\lsgfarbe #2\end{minipage}\par\vspace{-1.5mm}\else\linien{#1}\fi}
% Beliebiger Inhalt: \lsgoder{nur in der Lösung}{nur im Blatt}
\newcommand{\lsgoder}[2]{\ifloesung\leavevmode#1\else#2\fi}
% Text nur in der Lösung sichtbar, belegt im Blatt aber denselben Platz (ohne Text im PDF)
\newcommand{\lsgtext}[1]{\leavevmode\ifloesung{\lsgfarbe#1}\else\phantom{#1}\fi}
% Ankreuzfeld: \kreuz{1} = angekreuzt (nur in der Lösung), \kreuz{0} = leer
\newcommand{\kreuz}[1]{\leavevmode\ifloesung\ifnum#1=1 {\lsgfarbe$\boxtimes$}\else$\square$\fi\else$\square$\fi}
% Tabellenzeile mit fester Höhe für Lösungseinträge (2 Zeilen \small)
\newcommand{\zeile}{\rule[-6mm]{0pt}{9.5mm}}

% ---------- Verkettete Strukturen zeichnen ----------
% Objekt mit Kopf und Attributen:  \node[obj] (k1) {k1: Knoten \nodepart{second} daten = Hans\\ nachfolger};
\tikzset{
  obj/.style={draw,rectangle split,rectangle split parts=2,rectangle split part fill={orange!25,orange!8},
              rounded corners=2pt,font=\scriptsize,inner sep=2pt,align=left},
  qobj/.style={obj,rectangle split part fill={teal!25,teal!8}},
  aobj/.style={obj,rectangle split part fill={gray!30,gray!8}},
  pobj/.style={draw,rounded corners=2pt,fill=gray!8,font=\scriptsize,inner sep=2pt,align=left},
  klasse/.style={draw,rectangle split,rectangle split parts=3,font=\scriptsize,align=left,fill=white},
  zeiger/.style={thick,-{Stealth[length=2.2mm]}},
  erbt/.style={thick,-{Triangle[open,length=2.6mm,width=2.4mm]}},
  implementiert/.style={thick,densely dashed,-{Triangle[open,length=2.6mm,width=2.4mm]}},
  lsgzeiger/.style={zeiger,red!75!black},
  nullref/.style={font=\scriptsize\ttfamily,inner sep=1pt},
  schritt/.style={circle,draw,fill=white,inner sep=0.6pt,font=\scriptsize\bfseries},
  lsgschritt/.style={schritt,draw=red!75!black,text=red!75!black},
  % Sequenzdiagramm
  lebenslinie/.style={draw,rounded corners=1.5pt,fill=white,font=\scriptsize,inner sep=2pt},
  aufruf/.style={thick,-{Stealth[length=2mm]}},
  rueckgabe/.style={thick,densely dashed,-{Stealth[length=2mm,open]}},
}
% Lösungszeichnung: im Blatt unsichtbar, belegt aber denselben Platz
\ifloesung
  \tikzset{lsgbild/.style={red!75!black}}
\else
  \tikzset{lsgbild/.style={draw opacity=0,fill opacity=0,text opacity=0}}
\fi

% ---------- Java-Listings ----------
\lstset{language=Java,basicstyle=\ttfamily\small,keywordstyle=\bfseries,commentstyle=\color{gray!80!black},stringstyle=\color{black},showstringspaces=false,columns=flexible,keepspaces=true,tabsize=3,frame=single,framerule=0.4pt,rulecolor=\color{black},xleftmargin=2mm,framexleftmargin=1mm,aboveskip=2mm,belowskip=2mm,breaklines=true,escapechar=§,
 morekeywords={interface,implements,abstract,extends},
 literate={ä}{{\"a}}1 {ö}{{\"o}}1 {ü}{{\"u}}1 {Ä}{{\"A}}1 {Ö}{{\"O}}1 {Ü}{{\"U}}1 {ß}{{\ss}}1 {–}{{--}}1 {„}{{\glqq}}1 {“}{{\grqq}}1 {→}{{$\to$}}1}
\lstnewenvironment{java}[1][]{\lstset{#1}}{}

%  Übersetzen mit XeLaTeX, LuaLaTeX, pdfLaTeX oder Tectonic.
%
%  Blatt:   xelatex ab-1-2-interfaces.tex
%  Lösung:  xelatex ab-1-2-interfaces-lsg.tex
%           (Datei enthält nur: \def\mitloesung{}\documentclass[11pt,a4paper]{article}
\input{ab-vorlage}
\begin{document}
\abtitel{1.2}{Interfaces: ein Vertrag für alle Daten}

Am Bahnhof warten Taxis in einer Schlange: Wer zuerst ankommt, fährt zuerst. Die Stadt möchte dafür die \code{Warteschlange} aus 1.1 benutzen. Doch die kann bisher nur Patienten verwalten. Eine Kopie für Taxis (\code{TaxiKnoten}, \code{TaxiWarteschlange}) würde funktionieren, aber jede Verbesserung müsste man dann doppelt einbauen. \web{Station 1.2, Projekt Taxistand}

\begin{aufgabe}[Was muss die Warteschlange über ihre Daten wissen?]{1}
Untersuche im Projekt die Methoden der Warteschlange. Kreuze an, ob die Methode die Daten nur weiterreicht oder eine Methode der Daten aufruft. Notiere diese gegebenenfalls.\\[1mm]
\small\renewcommand{\arraystretch}{1.2}
\begin{tabularx}{\linewidth}{|l|c|c|X|}\hline
\textbf{Methode} & \textbf{reicht nur weiter} & \textbf{ruft eine Methode der Daten auf} & \textbf{welche?}\\\hline
\code{add(Patient p)} & \kreuz{1} & \kreuz{0} & \\\hline
\code{poll()}, \code{peek()} & \kreuz{1} & \kreuz{0} & \\\hline
\code{ausgeben()} & \kreuz{0} & \kreuz{1} & \lsgtext{\code{getName()}: der Name zum Ausgeben}\\\hline
\end{tabularx}\\[1.5mm]
\normalsize
\textbf{Folgerung:} Von ihren Daten muss die Warteschlange nur verlangen, dass sie einen \lsg[4.6cm]{Namen (Schlüsselwert) nennen} können. Genau das legen wir in einem Vertrag fest: im \textbf{Interface} \code{DatenElement}.
\end{aufgabe}

\begin{minipage}[t]{0.47\linewidth}
\vspace{0pt}
\begin{java}[basicstyle=\ttfamily\footnotesize]
public §\lsg[1.9cm]{interface}§ DatenElement {
   String gibSchluesselwert();
}

public class Taxi §\lsg[3.2cm]{implements DatenElement}§ {
   private int nummer;
   private String fahrer;
   ...
   public String gibSchluesselwert() {
      return §\lsg[2.6cm]{"Taxi " + nummer}§;
   }
}
\end{java}
\end{minipage}\hfill
\begin{minipage}[t]{0.505\linewidth}
\vspace{0pt}
\begin{merke}[Interface]
\small
\begin{itemize}
\item Ein Interface ist ein \lsg[1.7cm]{Vertrag}. Es legt nur \lsg[2.9cm]{Methodenköpfe} fest, keine Rümpfe.
\item Eine Klasse, die es implementiert (\code{implements}), muss \lsg[1.3cm]{alle} seine Methoden bereitstellen. Das prüft der \lsg[2cm]{Compiler}.
\item Ein Interface ist auch ein \lsg[1.9cm]{Datentyp}: Eine Variable vom Typ \code{DatenElement} kann Objekte jeder Klasse aufnehmen, die es implementiert. Objekte vom Interface selbst gibt es \lsg[1.2cm]{nicht}.
\end{itemize}
\end{merke}
\end{minipage}

\begin{aufgabe}[Klassendiagramm]{2}
\textbf{a)} Zeichne alle Beziehungen ein: Assoziationen mit Namen und Multiplizitäten, dazu \code{implements} als gestrichelten Pfeil mit leerem Dreieck zum Interface.\\
\textbf{b)} Umrande die Klassen der \textbf{Struktur} blau und die Klassen des \textbf{Inhalts} grün.\\[1mm]
\begin{tikzpicture}
\node[klasse,fill=teal!6,anchor=north west] (W) at (0,0)
  {\textbf{Warteschlange}\nodepart{second}- anfang: Knoten\\- ende: Knoten
   \nodepart{third}+ add(DatenElement d): boolean\\+ poll(): DatenElement\\+ ausgeben()};
\node[klasse,fill=orange!8,anchor=north west] (K) at (6.0,0)
  {\textbf{Knoten}\nodepart{second}- daten: DatenElement\\- nachfolger: Knoten
   \nodepart{third}+ getDaten(): DatenElement\\+ getNachfolger(): Knoten};
\node[klasse,fill=green!6,anchor=north,rectangle split parts=2] (D) at (13.1,0)
  {\textit{\guillemotleft interface\guillemotright}\\\textbf{DatenElement}\nodepart{second}+ gibSchluesselwert(): String};
\node[klasse,fill=gray!8,anchor=north] (P) at (11.2,-3.0)
  {\textbf{Patient}\nodepart{second}- name: String
   \nodepart{third}+ gibSchluesselwert(): String};
\node[klasse,fill=gray!8,anchor=north] (T) at (15.15,-3.0)
  {\textbf{Taxi}\nodepart{second}- nummer: int\\- fahrer: String
   \nodepart{third}+ gibSchluesselwert(): String};
\begin{scope}[font=\scriptsize]
  \draw[zeiger,lsgbild] ([yshift=-4mm]W.north east) -- ([yshift=-4mm]W.north east -| K.west)
     node[pos=0.3,above] {\lsgtext{anfang}} node[pos=0.86,below] {\lsgtext{0..1}};
  \draw[zeiger,lsgbild] ([yshift=-13mm]W.north east) -- ([yshift=-13mm]W.north east -| K.west)
     node[pos=0.3,below] {\lsgtext{ende}} node[pos=0.86,above] {\lsgtext{0..1}};
  \draw[zeiger,lsgbild] ([yshift=-5mm]K.north east) -- ([yshift=-5mm]K.north east -| D.west)
     node[pos=0.6,below] {\lsgtext{daten}} node[pos=0.88,above] {\lsgtext{1}};
  \draw[zeiger,lsgbild] ([xshift=-5mm]K.north east) -- ++(0,4mm) -| ([xshift=5mm]K.north west);
  \node[left] at ($(K.north west)+(4.5mm,3.2mm)$) {\lsgtext{nachfolger\quad 0..1}};
  \draw[implementiert,lsgbild] ([xshift=10mm]P.north) -- ([xshift=10mm]P.north |- D.south);
  \draw[implementiert,lsgbild] ([xshift=-9.5mm]T.north) -- ([xshift=-9.5mm]T.north |- D.south);
  \draw[lsgbild,blue!70!black,thick,rounded corners] ($(W.north west)+(-1.5mm,7.5mm)$) rectangle ($(K.east |- W.south)+(1.5mm,-1.5mm)$);
  % Inhalt: oben nur um das Interface, unten um Patient und Taxi (unter dem blauen Rahmen)
  \coordinate (gl) at ($(K.east)+(2.75mm,0)$);
  \coordinate (gt) at ($(D.north)+(0,2mm)$);
  \coordinate (gr) at ($(T.east)+(1.5mm,0)$);
  \coordinate (gb) at ($(T.south)+(0,-1.5mm)$);
  \coordinate (gp) at ($(P.west)+(-1.5mm,0)$);
  \coordinate (gm) at ($(W.south)+(0,-3mm)$);
  \draw[lsgbild,green!45!black,thick,rounded corners] (gl |- gt) -- (gr |- gt) -- (gr |- gb) -- (gp |- gb) -- (gp |- gm) -- (gl |- gm) -- cycle;
\end{scope}
\end{tikzpicture}\\[1mm]
\textbf{c)} Durch die Trennung von Struktur und Inhalt erledigt jede Klasse genau eine Aufgabe (hohe \lsg[2.1cm]{Kohäsion}). Struktur und Inhalt hängen nur über das Interface zusammen (lose \lsg[2.1cm]{Kopplung}). So kann die Warteschlange ohne jede Änderung \lsg[4.5cm]{jedes neue DatenElement} aufnehmen.
\end{aufgabe}

\newpage
\begin{aufgabe}[Projekt Taxistand, TODO 3a bis 3d]{3}
Baue das Projekt so um, dass Warteschlange und Knoten mit jedem \code{DatenElement} arbeiten. Ergänze die Tabelle und übertrage die Änderungen ins Projekt (Reihenfolge wie die TODOs).\\[1mm]
\small\renewcommand{\arraystretch}{1.25}
\begin{tabularx}{\linewidth}{|c|l|X|}\hline
 & \textbf{bisher} & \textbf{neu}\\\hline
3a & \code{DatenElement}: (leer) & \lsgtext{\code{String gibSchluesselwert();}}\\\hline
3b & \code{public class Patient \{} & \lsgtext{\code{public class Patient implements DatenElement \{}}\\
   & & \lsgtext{\code{+ public String gibSchluesselwert() \{ return name; \}}}\\\hline
3c & \code{private Patient daten;} & \lsgtext{\code{private DatenElement daten;}} \ (ebenso Konstruktor, \code{getDaten})\\\hline
3d & \code{public boolean add(Patient p)} & \lsgtext{\code{public boolean add(DatenElement daten)}}\\
   & \code{public Patient poll()} & \lsgtext{\code{public DatenElement poll()}} \ (ebenso \code{peek})\\
   & \code{aktuell.getDaten().getName()} & \lsgtext{\code{aktuell.getDaten().gibSchluesselwert()}}\\\hline
\end{tabularx}\\[1.5mm]
\normalsize
\textbf{b)} Nach dem Umbau von \code{Knoten} meldet der Compiler einen Fehler in \code{ausgeben()}. Erkläre!\\
\lsg[\linewidth]{Ein DatenElement hat nur die Methode gibSchluesselwert(), kein getName().}\\[1mm]
\lsg[\linewidth]{Die Warteschlange darf nur benutzen, was im Vertrag steht.}
\end{aufgabe}

\begin{aufgabe}[Projekt Taxistand, TODO 4]{4}
\textbf{a)} Ergänze den Code von \code{Taxi} oben auf Seite 1, übertrage ihn ins Projekt und kommentiere im Hauptprogramm die Taxi-Zeilen ein.\\
\textbf{b)} Lösche testweise die Methode \code{gibSchluesselwert()} aus \code{Taxi}. Notiere die Meldung des Compilers und erkläre sie mit dem Wort \textit{Vertrag}.\\
\lsg[\linewidth]{„Die Klasse Taxi muss noch folgende Methoden des Interfaces DatenElement implementieren: …“}\\[1mm]
\lsg[\linewidth]{Taxi verspricht mit implements, den Vertrag zu erfüllen, hält ihn aber nicht ein.}
\end{aufgabe}

\begin{aufgabe}[Datentypen]{5}
\code{taxistand} ist eine \code{Warteschlange}. Kreuze an, welche Zeilen der Compiler akzeptiert, und begründe kurz.\\[1mm]
\small\renewcommand{\arraystretch}{1.25}
\begin{tabularx}{\linewidth}{|l|c|c|X|}\hline
\textbf{Zeile} & \textbf{ok} & \textbf{Fehler} & \textbf{Begründung}\\\hline
\code{DatenElement d = new Taxi(5, "Lea");} & \kreuz{1} & \kreuz{0} & \lsgtext{Taxi implementiert DatenElement.}\\\hline
\code{DatenElement d = new DatenElement();} & \kreuz{0} & \kreuz{1} & \lsgtext{Von einem Interface gibt es keine Objekte.}\\\hline
\code{taxistand.add(new Patient("Ali", "Husten"));} & \kreuz{1} & \kreuz{0} & \lsgtext{Auch Patient ist ein DatenElement.}\\\hline
\code{Taxi t = taxistand.poll();} & \kreuz{0} & \kreuz{1} & \lsgtext{poll() liefert nur ein DatenElement.}\\\hline
\code{println(taxistand.poll().gibSchluesselwert());} & \kreuz{1} & \kreuz{0} & \lsgtext{Diese Methode steht im Vertrag.}\\\hline
\code{println(taxistand.poll().getFahrer());} & \kreuz{0} & \kreuz{1} & \lsgtext{getFahrer() steht nicht im Vertrag.}\\\hline
\end{tabularx}
\normalsize
\end{aufgabe}

\begin{hinweis}\small
\textbf{Es geht auch anders -- besprecht mündlich, warum man es so machen kann:}
\begin{itemize}
\item Jede Klasse ist ein \code{Object}. Der Knoten könnte also einfach \code{Object daten} speichern. Was ginge dann in \code{ausgeben()} nicht mehr?
\item In einer älteren Fassung stand im Vertrag nur \code{String toString();}. Warum erzwingt dieser Vertrag nichts? (Tipp: Was erbt jede Klasse von \code{Object}?)
\item Statt des Interfaces gäbe es eine Oberklasse \code{Daten}, von der Patient und Taxi erben. Was spricht dagegen, wenn \code{Patient} schon von \code{Person} erbt?
\item Du kennst schon ein Interface: \code{Wartezimmer} aus 1.1. Was legte es fest -- und wozu?
\end{itemize}
\end{hinweis}

\end{document}
;
%            füllt alle Lücken rot aus; siehe \lsg, \lsglinien, \lsgoder,
%            \lsgtext, lsgbild, \kreuz)
% =====================================================================
\usepackage[a4paper,left=16mm,right=16mm,top=23mm,bottom=13mm,headheight=12mm,headsep=3mm,footskip=6mm]{geometry}
\usepackage{iftex}
\ifPDFTeX
  \usepackage[T1]{fontenc}
  \usepackage[utf8]{inputenc}
  \usepackage{helvet}
  \usepackage[scaled=0.86]{beramono}
  \renewcommand{\familydefault}{\sfdefault}
\else
  \usepackage{fontspec}
  \setmainfont{texgyreheros}[Extension=.otf,UprightFont=*-regular,BoldFont=*-bold,ItalicFont=*-italic,BoldItalicFont=*-bolditalic]
  \setmonofont{DejaVuSansMono}[Extension=.ttf,UprightFont=*,BoldFont=*-Bold,ItalicFont=*-Oblique,BoldItalicFont=*-BoldOblique,Scale=0.86]
\fi
\usepackage[ngerman,shorthands=off]{babel}
\usepackage{amssymb}
\usepackage{xcolor,graphicx,tabularx,array,enumitem,fancyhdr,tikz,listings,multicol,calc,etoolbox,pifont}
\usepackage[most]{tcolorbox}
\usetikzlibrary{arrows.meta,positioning,calc,decorations.pathreplacing,shapes.multipart}
\usepackage[hidelinks]{hyperref}
\setlength{\parindent}{0pt}
\setlength{\parskip}{3pt}
\setlist{nosep,leftmargin=*}

% ---------- Lösungsschalter ----------
\newif\ifloesung
\ifdefined\mitloesung\loesungtrue\fi
\newcommand{\lsgfarbe}{\color{red!75!black}}

% ---------- Kopf- und Fußzeile (einheitliche Form) ----------
\newcommand{\lehrkraft}{Hau}
\newcommand{\themenbereich}{1. Rekursive Datenstrukturen}
\newcommand{\abnummer}{}
\pagestyle{fancy}\fancyhf{}
\renewcommand{\headrulewidth}{0.4pt}
\fancyhead[L]{\small Klasse 12\\Datum:}
\fancyhead[C]{\small Themenbereich:\\\themenbereich}
\fancyhead[R]{\small \lehrkraft\ifloesung\\{\lsgfarbe\bfseries Lösung}\fi}
\fancyfoot[C]{\scriptsize Informatik 12\,·\,Blatt \abnummer\ifloesung\ (Lösung)\fi\,·\,Seite \thepage}
\newcommand{\abtitel}[2]{\renewcommand{\abnummer}{#1}\begin{center}\Large\bfseries #1\quad\underline{#2}\end{center}\vspace{-2mm}}

% ---------- Boxen ----------
\newenvironment{aufgabe}[2][]{\begin{tcolorbox}[enhanced,breakable,colback=white,colframe=black,boxrule=0.7pt,arc=3mm,left=2mm,right=2mm,top=1mm,bottom=1.2mm,before skip=2mm,after skip=2mm]\textbf{Aufgabe #2)}\ifstrempty{#1}{}{ \textit{#1}}\par\vspace{0.3mm}}{\end{tcolorbox}}
\newtcolorbox{merke}[1][Merke]{enhanced,breakable,colback=green!5,colframe=green!40!black,boxrule=0.8pt,arc=2mm,left=2mm,right=2mm,top=1mm,bottom=1.2mm,before skip=2mm,after skip=2mm,title={#1},fonttitle=\bfseries,coltitle=white,toptitle=0.3mm,bottomtitle=0.3mm}
\newtcolorbox{hinweis}{enhanced,breakable,colback=yellow!12,colframe=orange!70!black,boxrule=0.6pt,arc=2mm,left=2mm,right=2mm,top=1mm,bottom=1mm,before skip=2mm,after skip=2mm}
\newcommand{\web}[1]{{\small\color{gray}\textit{Website: #1}}}

% ---------- Lücken, Linien, Karos ----------
\newcommand{\luecke}[1][3.5cm]{\rule[-0.5ex]{#1}{0.4pt}}
\newcommand{\linien}[1]{\par\vspace{1.5mm}\foreach \i in {1,...,#1}{\noindent\rule[0pt]{\linewidth}{0.3pt}\par\vspace{5.5mm}}\vspace{-3mm}}
\newcommand{\kariert}[1]{\par\vspace{1.5mm}\noindent\begin{tikzpicture}\draw[step=5mm,gray!55,thin] (0,0) grid (\linewidth-1pt,#1*5mm);\end{tikzpicture}\par}
\newcommand{\karobox}[2]{\begin{tikzpicture}\draw[step=5mm,gray!55,thin] (0,0) grid (#1,#2);\end{tikzpicture}}
\newcommand{\code}[1]{\texttt{#1}}
% Schritt-Nummern wie im Code und auf der Website: \nr{1} = ①
\newcommand{\nr}[1]{\ding{\numexpr171+#1\relax}}

% Lücke mit Lösung: \lsg[Mindestbreite]{Lösungstext}
%  Die Lösung steht (ohne eigene Höhe/Tiefe) auf derselben Linie wie im Blatt: Das Layout verschiebt sich nicht,
%  solange der Text nicht breiter als die Mindestbreite ist.
\newlength{\lsgbreite}
\newcommand{\lsg}[2][3.5cm]{\leavevmode\ifloesung\settowidth{\lsgbreite}{\,#2\,}\ifdim\lsgbreite<#1\setlength{\lsgbreite}{#1}\fi\rlap{\smash{\raisebox{0.8pt}{\makebox[\lsgbreite][l]{\lsgfarbe\,#2}}}}\luecke[\lsgbreite]\else\luecke[#1]\fi}
% Schreiblinien mit Lösung: \lsglinien{Anzahl}{Lösungstext} (gleiche Höhe wie \linien:
%  je Linie \baselineskip + 5,5 mm + \parskip, ausgemessen)
\newlength{\lsglinienhoehe}
\newcommand{\lsglinien}[2]{\ifloesung\par\vspace{1.5mm}\setlength{\lsglinienhoehe}{\dimexpr #1\dimexpr\baselineskip+5.5mm+\parskip\relax-16.25pt\relax}\noindent\begin{minipage}[t][\lsglinienhoehe][t]{\linewidth}\lsgfarbe #2\end{minipage}\par\vspace{-1.5mm}\else\linien{#1}\fi}
% Beliebiger Inhalt: \lsgoder{nur in der Lösung}{nur im Blatt}
\newcommand{\lsgoder}[2]{\ifloesung\leavevmode#1\else#2\fi}
% Text nur in der Lösung sichtbar, belegt im Blatt aber denselben Platz (ohne Text im PDF)
\newcommand{\lsgtext}[1]{\leavevmode\ifloesung{\lsgfarbe#1}\else\phantom{#1}\fi}
% Ankreuzfeld: \kreuz{1} = angekreuzt (nur in der Lösung), \kreuz{0} = leer
\newcommand{\kreuz}[1]{\leavevmode\ifloesung\ifnum#1=1 {\lsgfarbe$\boxtimes$}\else$\square$\fi\else$\square$\fi}
% Tabellenzeile mit fester Höhe für Lösungseinträge (2 Zeilen \small)
\newcommand{\zeile}{\rule[-6mm]{0pt}{9.5mm}}

% ---------- Verkettete Strukturen zeichnen ----------
% Objekt mit Kopf und Attributen:  \node[obj] (k1) {k1: Knoten \nodepart{second} daten = Hans\\ nachfolger};
\tikzset{
  obj/.style={draw,rectangle split,rectangle split parts=2,rectangle split part fill={orange!25,orange!8},
              rounded corners=2pt,font=\scriptsize,inner sep=2pt,align=left},
  qobj/.style={obj,rectangle split part fill={teal!25,teal!8}},
  aobj/.style={obj,rectangle split part fill={gray!30,gray!8}},
  pobj/.style={draw,rounded corners=2pt,fill=gray!8,font=\scriptsize,inner sep=2pt,align=left},
  klasse/.style={draw,rectangle split,rectangle split parts=3,font=\scriptsize,align=left,fill=white},
  zeiger/.style={thick,-{Stealth[length=2.2mm]}},
  erbt/.style={thick,-{Triangle[open,length=2.6mm,width=2.4mm]}},
  implementiert/.style={thick,densely dashed,-{Triangle[open,length=2.6mm,width=2.4mm]}},
  lsgzeiger/.style={zeiger,red!75!black},
  nullref/.style={font=\scriptsize\ttfamily,inner sep=1pt},
  schritt/.style={circle,draw,fill=white,inner sep=0.6pt,font=\scriptsize\bfseries},
  lsgschritt/.style={schritt,draw=red!75!black,text=red!75!black},
  % Sequenzdiagramm
  lebenslinie/.style={draw,rounded corners=1.5pt,fill=white,font=\scriptsize,inner sep=2pt},
  aufruf/.style={thick,-{Stealth[length=2mm]}},
  rueckgabe/.style={thick,densely dashed,-{Stealth[length=2mm,open]}},
}
% Lösungszeichnung: im Blatt unsichtbar, belegt aber denselben Platz
\ifloesung
  \tikzset{lsgbild/.style={red!75!black}}
\else
  \tikzset{lsgbild/.style={draw opacity=0,fill opacity=0,text opacity=0}}
\fi

% ---------- Java-Listings ----------
\lstset{language=Java,basicstyle=\ttfamily\small,keywordstyle=\bfseries,commentstyle=\color{gray!80!black},stringstyle=\color{black},showstringspaces=false,columns=flexible,keepspaces=true,tabsize=3,frame=single,framerule=0.4pt,rulecolor=\color{black},xleftmargin=2mm,framexleftmargin=1mm,aboveskip=2mm,belowskip=2mm,breaklines=true,escapechar=§,
 morekeywords={interface,implements,abstract,extends},
 literate={ä}{{\"a}}1 {ö}{{\"o}}1 {ü}{{\"u}}1 {Ä}{{\"A}}1 {Ö}{{\"O}}1 {Ü}{{\"U}}1 {ß}{{\ss}}1 {–}{{--}}1 {„}{{\glqq}}1 {“}{{\grqq}}1 {→}{{$\to$}}1}
\lstnewenvironment{java}[1][]{\lstset{#1}}{}

%  Übersetzen mit XeLaTeX, LuaLaTeX, pdfLaTeX oder Tectonic.
%
%  Blatt:   xelatex ab-1-2-interfaces.tex
%  Lösung:  xelatex ab-1-2-interfaces-lsg.tex
%           (Datei enthält nur: \def\mitloesung{}\documentclass[11pt,a4paper]{article}
% =====================================================================
%  Gemeinsame Vorlage der Arbeitsblätter "Rekursive Datenstrukturen"
%  – Informatik 12 (gleiches Layout wie informatik/13/formale-sprachen/
%  arbeitsblaetter/ab-vorlage.tex, dazu Listings und Objektdiagramme)
%  Einbinden in jedem Blatt direkt nach \documentclass: \input{ab-vorlage}
%  Übersetzen mit XeLaTeX, LuaLaTeX, pdfLaTeX oder Tectonic.
%
%  Blatt:   xelatex ab-1-2-interfaces.tex
%  Lösung:  xelatex ab-1-2-interfaces-lsg.tex
%           (Datei enthält nur: \def\mitloesung{}\input{ab-1-2-interfaces};
%            füllt alle Lücken rot aus; siehe \lsg, \lsglinien, \lsgoder,
%            \lsgtext, lsgbild, \kreuz)
% =====================================================================
\usepackage[a4paper,left=16mm,right=16mm,top=23mm,bottom=13mm,headheight=12mm,headsep=3mm,footskip=6mm]{geometry}
\usepackage{iftex}
\ifPDFTeX
  \usepackage[T1]{fontenc}
  \usepackage[utf8]{inputenc}
  \usepackage{helvet}
  \usepackage[scaled=0.86]{beramono}
  \renewcommand{\familydefault}{\sfdefault}
\else
  \usepackage{fontspec}
  \setmainfont{texgyreheros}[Extension=.otf,UprightFont=*-regular,BoldFont=*-bold,ItalicFont=*-italic,BoldItalicFont=*-bolditalic]
  \setmonofont{DejaVuSansMono}[Extension=.ttf,UprightFont=*,BoldFont=*-Bold,ItalicFont=*-Oblique,BoldItalicFont=*-BoldOblique,Scale=0.86]
\fi
\usepackage[ngerman,shorthands=off]{babel}
\usepackage{amssymb}
\usepackage{xcolor,graphicx,tabularx,array,enumitem,fancyhdr,tikz,listings,multicol,calc,etoolbox,pifont}
\usepackage[most]{tcolorbox}
\usetikzlibrary{arrows.meta,positioning,calc,decorations.pathreplacing,shapes.multipart}
\usepackage[hidelinks]{hyperref}
\setlength{\parindent}{0pt}
\setlength{\parskip}{3pt}
\setlist{nosep,leftmargin=*}

% ---------- Lösungsschalter ----------
\newif\ifloesung
\ifdefined\mitloesung\loesungtrue\fi
\newcommand{\lsgfarbe}{\color{red!75!black}}

% ---------- Kopf- und Fußzeile (einheitliche Form) ----------
\newcommand{\lehrkraft}{Hau}
\newcommand{\themenbereich}{1. Rekursive Datenstrukturen}
\newcommand{\abnummer}{}
\pagestyle{fancy}\fancyhf{}
\renewcommand{\headrulewidth}{0.4pt}
\fancyhead[L]{\small Klasse 12\\Datum:}
\fancyhead[C]{\small Themenbereich:\\\themenbereich}
\fancyhead[R]{\small \lehrkraft\ifloesung\\{\lsgfarbe\bfseries Lösung}\fi}
\fancyfoot[C]{\scriptsize Informatik 12\,·\,Blatt \abnummer\ifloesung\ (Lösung)\fi\,·\,Seite \thepage}
\newcommand{\abtitel}[2]{\renewcommand{\abnummer}{#1}\begin{center}\Large\bfseries #1\quad\underline{#2}\end{center}\vspace{-2mm}}

% ---------- Boxen ----------
\newenvironment{aufgabe}[2][]{\begin{tcolorbox}[enhanced,breakable,colback=white,colframe=black,boxrule=0.7pt,arc=3mm,left=2mm,right=2mm,top=1mm,bottom=1.2mm,before skip=2mm,after skip=2mm]\textbf{Aufgabe #2)}\ifstrempty{#1}{}{ \textit{#1}}\par\vspace{0.3mm}}{\end{tcolorbox}}
\newtcolorbox{merke}[1][Merke]{enhanced,breakable,colback=green!5,colframe=green!40!black,boxrule=0.8pt,arc=2mm,left=2mm,right=2mm,top=1mm,bottom=1.2mm,before skip=2mm,after skip=2mm,title={#1},fonttitle=\bfseries,coltitle=white,toptitle=0.3mm,bottomtitle=0.3mm}
\newtcolorbox{hinweis}{enhanced,breakable,colback=yellow!12,colframe=orange!70!black,boxrule=0.6pt,arc=2mm,left=2mm,right=2mm,top=1mm,bottom=1mm,before skip=2mm,after skip=2mm}
\newcommand{\web}[1]{{\small\color{gray}\textit{Website: #1}}}

% ---------- Lücken, Linien, Karos ----------
\newcommand{\luecke}[1][3.5cm]{\rule[-0.5ex]{#1}{0.4pt}}
\newcommand{\linien}[1]{\par\vspace{1.5mm}\foreach \i in {1,...,#1}{\noindent\rule[0pt]{\linewidth}{0.3pt}\par\vspace{5.5mm}}\vspace{-3mm}}
\newcommand{\kariert}[1]{\par\vspace{1.5mm}\noindent\begin{tikzpicture}\draw[step=5mm,gray!55,thin] (0,0) grid (\linewidth-1pt,#1*5mm);\end{tikzpicture}\par}
\newcommand{\karobox}[2]{\begin{tikzpicture}\draw[step=5mm,gray!55,thin] (0,0) grid (#1,#2);\end{tikzpicture}}
\newcommand{\code}[1]{\texttt{#1}}
% Schritt-Nummern wie im Code und auf der Website: \nr{1} = ①
\newcommand{\nr}[1]{\ding{\numexpr171+#1\relax}}

% Lücke mit Lösung: \lsg[Mindestbreite]{Lösungstext}
%  Die Lösung steht (ohne eigene Höhe/Tiefe) auf derselben Linie wie im Blatt: Das Layout verschiebt sich nicht,
%  solange der Text nicht breiter als die Mindestbreite ist.
\newlength{\lsgbreite}
\newcommand{\lsg}[2][3.5cm]{\leavevmode\ifloesung\settowidth{\lsgbreite}{\,#2\,}\ifdim\lsgbreite<#1\setlength{\lsgbreite}{#1}\fi\rlap{\smash{\raisebox{0.8pt}{\makebox[\lsgbreite][l]{\lsgfarbe\,#2}}}}\luecke[\lsgbreite]\else\luecke[#1]\fi}
% Schreiblinien mit Lösung: \lsglinien{Anzahl}{Lösungstext} (gleiche Höhe wie \linien:
%  je Linie \baselineskip + 5,5 mm + \parskip, ausgemessen)
\newlength{\lsglinienhoehe}
\newcommand{\lsglinien}[2]{\ifloesung\par\vspace{1.5mm}\setlength{\lsglinienhoehe}{\dimexpr #1\dimexpr\baselineskip+5.5mm+\parskip\relax-16.25pt\relax}\noindent\begin{minipage}[t][\lsglinienhoehe][t]{\linewidth}\lsgfarbe #2\end{minipage}\par\vspace{-1.5mm}\else\linien{#1}\fi}
% Beliebiger Inhalt: \lsgoder{nur in der Lösung}{nur im Blatt}
\newcommand{\lsgoder}[2]{\ifloesung\leavevmode#1\else#2\fi}
% Text nur in der Lösung sichtbar, belegt im Blatt aber denselben Platz (ohne Text im PDF)
\newcommand{\lsgtext}[1]{\leavevmode\ifloesung{\lsgfarbe#1}\else\phantom{#1}\fi}
% Ankreuzfeld: \kreuz{1} = angekreuzt (nur in der Lösung), \kreuz{0} = leer
\newcommand{\kreuz}[1]{\leavevmode\ifloesung\ifnum#1=1 {\lsgfarbe$\boxtimes$}\else$\square$\fi\else$\square$\fi}
% Tabellenzeile mit fester Höhe für Lösungseinträge (2 Zeilen \small)
\newcommand{\zeile}{\rule[-6mm]{0pt}{9.5mm}}

% ---------- Verkettete Strukturen zeichnen ----------
% Objekt mit Kopf und Attributen:  \node[obj] (k1) {k1: Knoten \nodepart{second} daten = Hans\\ nachfolger};
\tikzset{
  obj/.style={draw,rectangle split,rectangle split parts=2,rectangle split part fill={orange!25,orange!8},
              rounded corners=2pt,font=\scriptsize,inner sep=2pt,align=left},
  qobj/.style={obj,rectangle split part fill={teal!25,teal!8}},
  aobj/.style={obj,rectangle split part fill={gray!30,gray!8}},
  pobj/.style={draw,rounded corners=2pt,fill=gray!8,font=\scriptsize,inner sep=2pt,align=left},
  klasse/.style={draw,rectangle split,rectangle split parts=3,font=\scriptsize,align=left,fill=white},
  zeiger/.style={thick,-{Stealth[length=2.2mm]}},
  erbt/.style={thick,-{Triangle[open,length=2.6mm,width=2.4mm]}},
  implementiert/.style={thick,densely dashed,-{Triangle[open,length=2.6mm,width=2.4mm]}},
  lsgzeiger/.style={zeiger,red!75!black},
  nullref/.style={font=\scriptsize\ttfamily,inner sep=1pt},
  schritt/.style={circle,draw,fill=white,inner sep=0.6pt,font=\scriptsize\bfseries},
  lsgschritt/.style={schritt,draw=red!75!black,text=red!75!black},
  % Sequenzdiagramm
  lebenslinie/.style={draw,rounded corners=1.5pt,fill=white,font=\scriptsize,inner sep=2pt},
  aufruf/.style={thick,-{Stealth[length=2mm]}},
  rueckgabe/.style={thick,densely dashed,-{Stealth[length=2mm,open]}},
}
% Lösungszeichnung: im Blatt unsichtbar, belegt aber denselben Platz
\ifloesung
  \tikzset{lsgbild/.style={red!75!black}}
\else
  \tikzset{lsgbild/.style={draw opacity=0,fill opacity=0,text opacity=0}}
\fi

% ---------- Java-Listings ----------
\lstset{language=Java,basicstyle=\ttfamily\small,keywordstyle=\bfseries,commentstyle=\color{gray!80!black},stringstyle=\color{black},showstringspaces=false,columns=flexible,keepspaces=true,tabsize=3,frame=single,framerule=0.4pt,rulecolor=\color{black},xleftmargin=2mm,framexleftmargin=1mm,aboveskip=2mm,belowskip=2mm,breaklines=true,escapechar=§,
 morekeywords={interface,implements,abstract,extends},
 literate={ä}{{\"a}}1 {ö}{{\"o}}1 {ü}{{\"u}}1 {Ä}{{\"A}}1 {Ö}{{\"O}}1 {Ü}{{\"U}}1 {ß}{{\ss}}1 {–}{{--}}1 {„}{{\glqq}}1 {“}{{\grqq}}1 {→}{{$\to$}}1}
\lstnewenvironment{java}[1][]{\lstset{#1}}{}

\begin{document}
\abtitel{1.2}{Interfaces: ein Vertrag für alle Daten}

Am Bahnhof warten Taxis in einer Schlange: Wer zuerst ankommt, fährt zuerst. Die Stadt möchte dafür die \code{Warteschlange} aus 1.1 benutzen. Doch die kann bisher nur Patienten verwalten. Eine Kopie für Taxis (\code{TaxiKnoten}, \code{TaxiWarteschlange}) würde funktionieren, aber jede Verbesserung müsste man dann doppelt einbauen. \web{Station 1.2, Projekt Taxistand}

\begin{aufgabe}[Was muss die Warteschlange über ihre Daten wissen?]{1}
Untersuche im Projekt die Methoden der Warteschlange. Kreuze an, ob die Methode die Daten nur weiterreicht oder eine Methode der Daten aufruft. Notiere diese gegebenenfalls.\\[1mm]
\small\renewcommand{\arraystretch}{1.2}
\begin{tabularx}{\linewidth}{|l|c|c|X|}\hline
\textbf{Methode} & \textbf{reicht nur weiter} & \textbf{ruft eine Methode der Daten auf} & \textbf{welche?}\\\hline
\code{add(Patient p)} & \kreuz{1} & \kreuz{0} & \\\hline
\code{poll()}, \code{peek()} & \kreuz{1} & \kreuz{0} & \\\hline
\code{ausgeben()} & \kreuz{0} & \kreuz{1} & \lsgtext{\code{getName()}: der Name zum Ausgeben}\\\hline
\end{tabularx}\\[1.5mm]
\normalsize
\textbf{Folgerung:} Von ihren Daten muss die Warteschlange nur verlangen, dass sie einen \lsg[4.6cm]{Namen (Schlüsselwert) nennen} können. Genau das legen wir in einem Vertrag fest: im \textbf{Interface} \code{DatenElement}.
\end{aufgabe}

\begin{minipage}[t]{0.47\linewidth}
\vspace{0pt}
\begin{java}[basicstyle=\ttfamily\footnotesize]
public §\lsg[1.9cm]{interface}§ DatenElement {
   String gibSchluesselwert();
}

public class Taxi §\lsg[3.2cm]{implements DatenElement}§ {
   private int nummer;
   private String fahrer;
   ...
   public String gibSchluesselwert() {
      return §\lsg[2.6cm]{"Taxi " + nummer}§;
   }
}
\end{java}
\end{minipage}\hfill
\begin{minipage}[t]{0.505\linewidth}
\vspace{0pt}
\begin{merke}[Interface]
\small
\begin{itemize}
\item Ein Interface ist ein \lsg[1.7cm]{Vertrag}. Es legt nur \lsg[2.9cm]{Methodenköpfe} fest, keine Rümpfe.
\item Eine Klasse, die es implementiert (\code{implements}), muss \lsg[1.3cm]{alle} seine Methoden bereitstellen. Das prüft der \lsg[2cm]{Compiler}.
\item Ein Interface ist auch ein \lsg[1.9cm]{Datentyp}: Eine Variable vom Typ \code{DatenElement} kann Objekte jeder Klasse aufnehmen, die es implementiert. Objekte vom Interface selbst gibt es \lsg[1.2cm]{nicht}.
\end{itemize}
\end{merke}
\end{minipage}

\begin{aufgabe}[Klassendiagramm]{2}
\textbf{a)} Zeichne alle Beziehungen ein: Assoziationen mit Namen und Multiplizitäten, dazu \code{implements} als gestrichelten Pfeil mit leerem Dreieck zum Interface.\\
\textbf{b)} Umrande die Klassen der \textbf{Struktur} blau und die Klassen des \textbf{Inhalts} grün.\\[1mm]
\begin{tikzpicture}
\node[klasse,fill=teal!6,anchor=north west] (W) at (0,0)
  {\textbf{Warteschlange}\nodepart{second}- anfang: Knoten\\- ende: Knoten
   \nodepart{third}+ add(DatenElement d): boolean\\+ poll(): DatenElement\\+ ausgeben()};
\node[klasse,fill=orange!8,anchor=north west] (K) at (6.0,0)
  {\textbf{Knoten}\nodepart{second}- daten: DatenElement\\- nachfolger: Knoten
   \nodepart{third}+ getDaten(): DatenElement\\+ getNachfolger(): Knoten};
\node[klasse,fill=green!6,anchor=north,rectangle split parts=2] (D) at (13.1,0)
  {\textit{\guillemotleft interface\guillemotright}\\\textbf{DatenElement}\nodepart{second}+ gibSchluesselwert(): String};
\node[klasse,fill=gray!8,anchor=north] (P) at (11.2,-3.0)
  {\textbf{Patient}\nodepart{second}- name: String
   \nodepart{third}+ gibSchluesselwert(): String};
\node[klasse,fill=gray!8,anchor=north] (T) at (15.15,-3.0)
  {\textbf{Taxi}\nodepart{second}- nummer: int\\- fahrer: String
   \nodepart{third}+ gibSchluesselwert(): String};
\begin{scope}[font=\scriptsize]
  \draw[zeiger,lsgbild] ([yshift=-4mm]W.north east) -- ([yshift=-4mm]W.north east -| K.west)
     node[pos=0.3,above] {\lsgtext{anfang}} node[pos=0.86,below] {\lsgtext{0..1}};
  \draw[zeiger,lsgbild] ([yshift=-13mm]W.north east) -- ([yshift=-13mm]W.north east -| K.west)
     node[pos=0.3,below] {\lsgtext{ende}} node[pos=0.86,above] {\lsgtext{0..1}};
  \draw[zeiger,lsgbild] ([yshift=-5mm]K.north east) -- ([yshift=-5mm]K.north east -| D.west)
     node[pos=0.6,below] {\lsgtext{daten}} node[pos=0.88,above] {\lsgtext{1}};
  \draw[zeiger,lsgbild] ([xshift=-5mm]K.north east) -- ++(0,4mm) -| ([xshift=5mm]K.north west);
  \node[left] at ($(K.north west)+(4.5mm,3.2mm)$) {\lsgtext{nachfolger\quad 0..1}};
  \draw[implementiert,lsgbild] ([xshift=10mm]P.north) -- ([xshift=10mm]P.north |- D.south);
  \draw[implementiert,lsgbild] ([xshift=-9.5mm]T.north) -- ([xshift=-9.5mm]T.north |- D.south);
  \draw[lsgbild,blue!70!black,thick,rounded corners] ($(W.north west)+(-1.5mm,7.5mm)$) rectangle ($(K.east |- W.south)+(1.5mm,-1.5mm)$);
  % Inhalt: oben nur um das Interface, unten um Patient und Taxi (unter dem blauen Rahmen)
  \coordinate (gl) at ($(K.east)+(2.75mm,0)$);
  \coordinate (gt) at ($(D.north)+(0,2mm)$);
  \coordinate (gr) at ($(T.east)+(1.5mm,0)$);
  \coordinate (gb) at ($(T.south)+(0,-1.5mm)$);
  \coordinate (gp) at ($(P.west)+(-1.5mm,0)$);
  \coordinate (gm) at ($(W.south)+(0,-3mm)$);
  \draw[lsgbild,green!45!black,thick,rounded corners] (gl |- gt) -- (gr |- gt) -- (gr |- gb) -- (gp |- gb) -- (gp |- gm) -- (gl |- gm) -- cycle;
\end{scope}
\end{tikzpicture}\\[1mm]
\textbf{c)} Durch die Trennung von Struktur und Inhalt erledigt jede Klasse genau eine Aufgabe (hohe \lsg[2.1cm]{Kohäsion}). Struktur und Inhalt hängen nur über das Interface zusammen (lose \lsg[2.1cm]{Kopplung}). So kann die Warteschlange ohne jede Änderung \lsg[4.5cm]{jedes neue DatenElement} aufnehmen.
\end{aufgabe}

\newpage
\begin{aufgabe}[Projekt Taxistand, TODO 3a bis 3d]{3}
Baue das Projekt so um, dass Warteschlange und Knoten mit jedem \code{DatenElement} arbeiten. Ergänze die Tabelle und übertrage die Änderungen ins Projekt (Reihenfolge wie die TODOs).\\[1mm]
\small\renewcommand{\arraystretch}{1.25}
\begin{tabularx}{\linewidth}{|c|l|X|}\hline
 & \textbf{bisher} & \textbf{neu}\\\hline
3a & \code{DatenElement}: (leer) & \lsgtext{\code{String gibSchluesselwert();}}\\\hline
3b & \code{public class Patient \{} & \lsgtext{\code{public class Patient implements DatenElement \{}}\\
   & & \lsgtext{\code{+ public String gibSchluesselwert() \{ return name; \}}}\\\hline
3c & \code{private Patient daten;} & \lsgtext{\code{private DatenElement daten;}} \ (ebenso Konstruktor, \code{getDaten})\\\hline
3d & \code{public boolean add(Patient p)} & \lsgtext{\code{public boolean add(DatenElement daten)}}\\
   & \code{public Patient poll()} & \lsgtext{\code{public DatenElement poll()}} \ (ebenso \code{peek})\\
   & \code{aktuell.getDaten().getName()} & \lsgtext{\code{aktuell.getDaten().gibSchluesselwert()}}\\\hline
\end{tabularx}\\[1.5mm]
\normalsize
\textbf{b)} Nach dem Umbau von \code{Knoten} meldet der Compiler einen Fehler in \code{ausgeben()}. Erkläre!\\
\lsg[\linewidth]{Ein DatenElement hat nur die Methode gibSchluesselwert(), kein getName().}\\[1mm]
\lsg[\linewidth]{Die Warteschlange darf nur benutzen, was im Vertrag steht.}
\end{aufgabe}

\begin{aufgabe}[Projekt Taxistand, TODO 4]{4}
\textbf{a)} Ergänze den Code von \code{Taxi} oben auf Seite 1, übertrage ihn ins Projekt und kommentiere im Hauptprogramm die Taxi-Zeilen ein.\\
\textbf{b)} Lösche testweise die Methode \code{gibSchluesselwert()} aus \code{Taxi}. Notiere die Meldung des Compilers und erkläre sie mit dem Wort \textit{Vertrag}.\\
\lsg[\linewidth]{„Die Klasse Taxi muss noch folgende Methoden des Interfaces DatenElement implementieren: …“}\\[1mm]
\lsg[\linewidth]{Taxi verspricht mit implements, den Vertrag zu erfüllen, hält ihn aber nicht ein.}
\end{aufgabe}

\begin{aufgabe}[Datentypen]{5}
\code{taxistand} ist eine \code{Warteschlange}. Kreuze an, welche Zeilen der Compiler akzeptiert, und begründe kurz.\\[1mm]
\small\renewcommand{\arraystretch}{1.25}
\begin{tabularx}{\linewidth}{|l|c|c|X|}\hline
\textbf{Zeile} & \textbf{ok} & \textbf{Fehler} & \textbf{Begründung}\\\hline
\code{DatenElement d = new Taxi(5, "Lea");} & \kreuz{1} & \kreuz{0} & \lsgtext{Taxi implementiert DatenElement.}\\\hline
\code{DatenElement d = new DatenElement();} & \kreuz{0} & \kreuz{1} & \lsgtext{Von einem Interface gibt es keine Objekte.}\\\hline
\code{taxistand.add(new Patient("Ali", "Husten"));} & \kreuz{1} & \kreuz{0} & \lsgtext{Auch Patient ist ein DatenElement.}\\\hline
\code{Taxi t = taxistand.poll();} & \kreuz{0} & \kreuz{1} & \lsgtext{poll() liefert nur ein DatenElement.}\\\hline
\code{println(taxistand.poll().gibSchluesselwert());} & \kreuz{1} & \kreuz{0} & \lsgtext{Diese Methode steht im Vertrag.}\\\hline
\code{println(taxistand.poll().getFahrer());} & \kreuz{0} & \kreuz{1} & \lsgtext{getFahrer() steht nicht im Vertrag.}\\\hline
\end{tabularx}
\normalsize
\end{aufgabe}

\begin{hinweis}\small
\textbf{Es geht auch anders -- besprecht mündlich, warum man es so machen kann:}
\begin{itemize}
\item Jede Klasse ist ein \code{Object}. Der Knoten könnte also einfach \code{Object daten} speichern. Was ginge dann in \code{ausgeben()} nicht mehr?
\item In einer älteren Fassung stand im Vertrag nur \code{String toString();}. Warum erzwingt dieser Vertrag nichts? (Tipp: Was erbt jede Klasse von \code{Object}?)
\item Statt des Interfaces gäbe es eine Oberklasse \code{Daten}, von der Patient und Taxi erben. Was spricht dagegen, wenn \code{Patient} schon von \code{Person} erbt?
\item Du kennst schon ein Interface: \code{Wartezimmer} aus 1.1. Was legte es fest -- und wozu?
\end{itemize}
\end{hinweis}

\end{document}
;
%            füllt alle Lücken rot aus; siehe \lsg, \lsglinien, \lsgoder,
%            \lsgtext, lsgbild, \kreuz)
% =====================================================================
\usepackage[a4paper,left=16mm,right=16mm,top=23mm,bottom=13mm,headheight=12mm,headsep=3mm,footskip=6mm]{geometry}
\usepackage{iftex}
\ifPDFTeX
  \usepackage[T1]{fontenc}
  \usepackage[utf8]{inputenc}
  \usepackage{helvet}
  \usepackage[scaled=0.86]{beramono}
  \renewcommand{\familydefault}{\sfdefault}
\else
  \usepackage{fontspec}
  \setmainfont{texgyreheros}[Extension=.otf,UprightFont=*-regular,BoldFont=*-bold,ItalicFont=*-italic,BoldItalicFont=*-bolditalic]
  \setmonofont{DejaVuSansMono}[Extension=.ttf,UprightFont=*,BoldFont=*-Bold,ItalicFont=*-Oblique,BoldItalicFont=*-BoldOblique,Scale=0.86]
\fi
\usepackage[ngerman,shorthands=off]{babel}
\usepackage{amssymb}
\usepackage{xcolor,graphicx,tabularx,array,enumitem,fancyhdr,tikz,listings,multicol,calc,etoolbox,pifont}
\usepackage[most]{tcolorbox}
\usetikzlibrary{arrows.meta,positioning,calc,decorations.pathreplacing,shapes.multipart}
\usepackage[hidelinks]{hyperref}
\setlength{\parindent}{0pt}
\setlength{\parskip}{3pt}
\setlist{nosep,leftmargin=*}

% ---------- Lösungsschalter ----------
\newif\ifloesung
\ifdefined\mitloesung\loesungtrue\fi
\newcommand{\lsgfarbe}{\color{red!75!black}}

% ---------- Kopf- und Fußzeile (einheitliche Form) ----------
\newcommand{\lehrkraft}{Hau}
\newcommand{\themenbereich}{1. Rekursive Datenstrukturen}
\newcommand{\abnummer}{}
\pagestyle{fancy}\fancyhf{}
\renewcommand{\headrulewidth}{0.4pt}
\fancyhead[L]{\small Klasse 12\\Datum:}
\fancyhead[C]{\small Themenbereich:\\\themenbereich}
\fancyhead[R]{\small \lehrkraft\ifloesung\\{\lsgfarbe\bfseries Lösung}\fi}
\fancyfoot[C]{\scriptsize Informatik 12\,·\,Blatt \abnummer\ifloesung\ (Lösung)\fi\,·\,Seite \thepage}
\newcommand{\abtitel}[2]{\renewcommand{\abnummer}{#1}\begin{center}\Large\bfseries #1\quad\underline{#2}\end{center}\vspace{-2mm}}

% ---------- Boxen ----------
\newenvironment{aufgabe}[2][]{\begin{tcolorbox}[enhanced,breakable,colback=white,colframe=black,boxrule=0.7pt,arc=3mm,left=2mm,right=2mm,top=1mm,bottom=1.2mm,before skip=2mm,after skip=2mm]\textbf{Aufgabe #2)}\ifstrempty{#1}{}{ \textit{#1}}\par\vspace{0.3mm}}{\end{tcolorbox}}
\newtcolorbox{merke}[1][Merke]{enhanced,breakable,colback=green!5,colframe=green!40!black,boxrule=0.8pt,arc=2mm,left=2mm,right=2mm,top=1mm,bottom=1.2mm,before skip=2mm,after skip=2mm,title={#1},fonttitle=\bfseries,coltitle=white,toptitle=0.3mm,bottomtitle=0.3mm}
\newtcolorbox{hinweis}{enhanced,breakable,colback=yellow!12,colframe=orange!70!black,boxrule=0.6pt,arc=2mm,left=2mm,right=2mm,top=1mm,bottom=1mm,before skip=2mm,after skip=2mm}
\newcommand{\web}[1]{{\small\color{gray}\textit{Website: #1}}}

% ---------- Lücken, Linien, Karos ----------
\newcommand{\luecke}[1][3.5cm]{\rule[-0.5ex]{#1}{0.4pt}}
\newcommand{\linien}[1]{\par\vspace{1.5mm}\foreach \i in {1,...,#1}{\noindent\rule[0pt]{\linewidth}{0.3pt}\par\vspace{5.5mm}}\vspace{-3mm}}
\newcommand{\kariert}[1]{\par\vspace{1.5mm}\noindent\begin{tikzpicture}\draw[step=5mm,gray!55,thin] (0,0) grid (\linewidth-1pt,#1*5mm);\end{tikzpicture}\par}
\newcommand{\karobox}[2]{\begin{tikzpicture}\draw[step=5mm,gray!55,thin] (0,0) grid (#1,#2);\end{tikzpicture}}
\newcommand{\code}[1]{\texttt{#1}}
% Schritt-Nummern wie im Code und auf der Website: \nr{1} = ①
\newcommand{\nr}[1]{\ding{\numexpr171+#1\relax}}

% Lücke mit Lösung: \lsg[Mindestbreite]{Lösungstext}
%  Die Lösung steht (ohne eigene Höhe/Tiefe) auf derselben Linie wie im Blatt: Das Layout verschiebt sich nicht,
%  solange der Text nicht breiter als die Mindestbreite ist.
\newlength{\lsgbreite}
\newcommand{\lsg}[2][3.5cm]{\leavevmode\ifloesung\settowidth{\lsgbreite}{\,#2\,}\ifdim\lsgbreite<#1\setlength{\lsgbreite}{#1}\fi\rlap{\smash{\raisebox{0.8pt}{\makebox[\lsgbreite][l]{\lsgfarbe\,#2}}}}\luecke[\lsgbreite]\else\luecke[#1]\fi}
% Schreiblinien mit Lösung: \lsglinien{Anzahl}{Lösungstext} (gleiche Höhe wie \linien:
%  je Linie \baselineskip + 5,5 mm + \parskip, ausgemessen)
\newlength{\lsglinienhoehe}
\newcommand{\lsglinien}[2]{\ifloesung\par\vspace{1.5mm}\setlength{\lsglinienhoehe}{\dimexpr #1\dimexpr\baselineskip+5.5mm+\parskip\relax-16.25pt\relax}\noindent\begin{minipage}[t][\lsglinienhoehe][t]{\linewidth}\lsgfarbe #2\end{minipage}\par\vspace{-1.5mm}\else\linien{#1}\fi}
% Beliebiger Inhalt: \lsgoder{nur in der Lösung}{nur im Blatt}
\newcommand{\lsgoder}[2]{\ifloesung\leavevmode#1\else#2\fi}
% Text nur in der Lösung sichtbar, belegt im Blatt aber denselben Platz (ohne Text im PDF)
\newcommand{\lsgtext}[1]{\leavevmode\ifloesung{\lsgfarbe#1}\else\phantom{#1}\fi}
% Ankreuzfeld: \kreuz{1} = angekreuzt (nur in der Lösung), \kreuz{0} = leer
\newcommand{\kreuz}[1]{\leavevmode\ifloesung\ifnum#1=1 {\lsgfarbe$\boxtimes$}\else$\square$\fi\else$\square$\fi}
% Tabellenzeile mit fester Höhe für Lösungseinträge (2 Zeilen \small)
\newcommand{\zeile}{\rule[-6mm]{0pt}{9.5mm}}

% ---------- Verkettete Strukturen zeichnen ----------
% Objekt mit Kopf und Attributen:  \node[obj] (k1) {k1: Knoten \nodepart{second} daten = Hans\\ nachfolger};
\tikzset{
  obj/.style={draw,rectangle split,rectangle split parts=2,rectangle split part fill={orange!25,orange!8},
              rounded corners=2pt,font=\scriptsize,inner sep=2pt,align=left},
  qobj/.style={obj,rectangle split part fill={teal!25,teal!8}},
  aobj/.style={obj,rectangle split part fill={gray!30,gray!8}},
  pobj/.style={draw,rounded corners=2pt,fill=gray!8,font=\scriptsize,inner sep=2pt,align=left},
  klasse/.style={draw,rectangle split,rectangle split parts=3,font=\scriptsize,align=left,fill=white},
  zeiger/.style={thick,-{Stealth[length=2.2mm]}},
  erbt/.style={thick,-{Triangle[open,length=2.6mm,width=2.4mm]}},
  implementiert/.style={thick,densely dashed,-{Triangle[open,length=2.6mm,width=2.4mm]}},
  lsgzeiger/.style={zeiger,red!75!black},
  nullref/.style={font=\scriptsize\ttfamily,inner sep=1pt},
  schritt/.style={circle,draw,fill=white,inner sep=0.6pt,font=\scriptsize\bfseries},
  lsgschritt/.style={schritt,draw=red!75!black,text=red!75!black},
  % Sequenzdiagramm
  lebenslinie/.style={draw,rounded corners=1.5pt,fill=white,font=\scriptsize,inner sep=2pt},
  aufruf/.style={thick,-{Stealth[length=2mm]}},
  rueckgabe/.style={thick,densely dashed,-{Stealth[length=2mm,open]}},
}
% Lösungszeichnung: im Blatt unsichtbar, belegt aber denselben Platz
\ifloesung
  \tikzset{lsgbild/.style={red!75!black}}
\else
  \tikzset{lsgbild/.style={draw opacity=0,fill opacity=0,text opacity=0}}
\fi

% ---------- Java-Listings ----------
\lstset{language=Java,basicstyle=\ttfamily\small,keywordstyle=\bfseries,commentstyle=\color{gray!80!black},stringstyle=\color{black},showstringspaces=false,columns=flexible,keepspaces=true,tabsize=3,frame=single,framerule=0.4pt,rulecolor=\color{black},xleftmargin=2mm,framexleftmargin=1mm,aboveskip=2mm,belowskip=2mm,breaklines=true,escapechar=§,
 morekeywords={interface,implements,abstract,extends},
 literate={ä}{{\"a}}1 {ö}{{\"o}}1 {ü}{{\"u}}1 {Ä}{{\"A}}1 {Ö}{{\"O}}1 {Ü}{{\"U}}1 {ß}{{\ss}}1 {–}{{--}}1 {„}{{\glqq}}1 {“}{{\grqq}}1 {→}{{$\to$}}1}
\lstnewenvironment{java}[1][]{\lstset{#1}}{}

\begin{document}
\abtitel{1.1}{Die Warteschlange (Queue)}

In der Arztpraxis wird in der Reihenfolge der Ankunft aufgerufen. Das \textbf{Projekt Praxis} \web{Station 1.1} verwaltet das Wartezimmer zunächst mit einem \textbf{Array}. Jedes Wartezimmer kann \code{add(Patient p)} (hinten anstellen), \code{poll()} (vorne aufrufen) und \code{ausgeben()}.

\begin{aufgabe}[Projekt Praxis \quad\web{Wartezimmer-Simulator}]{1}
Starte das Hauptprogramm nacheinander mit Variante 1, 2 und 3. Notiere Idee und Problem.\\[1mm]
\small\renewcommand{\arraystretch}{1.15}
\begin{tabularx}{\linewidth}{|p{2.9cm}|X|X|}\hline
\textbf{Variante} & \textbf{Idee} & \textbf{Was geht schief?}\\\hline
1 \code{Suchen}\zeile & \lsgoder{\lsgfarbe add: erster freier Platz, poll: erster belegter Platz}{} & \lsgoder{\lsgfarbe Neue rücken vorne in freie Plätze und kommen vor Resi und Bene dran.}{}\\\hline
2 \code{Nummer}\zeile & \lsgoder{\lsgfarbe Eine Nummer merkt sich, welcher Platz als nächster dran ist.}{} & \lsgoder{\lsgfarbe Vorne Belegte kommen nie dran, die Nummer läuft übers Array hinaus.}{}\\\hline
3 \code{Aufrücken}\zeile & \lsgoder{\lsgfarbe Nach jedem Aufruf rücken alle einen Platz vor.}{} & \lsgoder{\lsgfarbe Reihenfolge stimmt, aber bei jedem Aufruf werden alle verschoben.}{}\\\hline
\end{tabularx}\\[1.5mm]
Was passiert bei \textbf{allen drei} Varianten mit Mia? \lsg[9.2cm]{Sie wird abgewiesen: Das Array ist voll (5 Plätze).}
\end{aufgabe}

\begin{merke}[Probleme mit Arrays]
Ein Array hat eine \lsg[2.6cm]{feste Größe} und belegt schon beim Erzeugen \lsg[2.2cm]{Speicher} für alle Plätze, auch für leere (\code{null}). Zum Vergrößern müsste man \lsg[3.2cm]{ein neues Array} erzeugen und alles \lsg[2.2cm]{kopieren}. Damit die Reihenfolge stimmt, müssen beim Herausnehmen alle \lsg[2.4cm]{aufrücken}.\\[0.5mm]
\textbf{Idee:} Jedes Element kennt zusätzlich \lsg[3.4cm]{seinen Nachfolger}. So entsteht eine \textbf{Kette} beliebiger Länge.
\end{merke}

\begin{aufgabe}{2}
\textbf{a)} Resi, Bene und Gustav warten. Zeichne alle Verweise als Pfeile ein (oder \code{null}). Markiere, wo \code{poll()} herausnimmt und wo \code{add()} anhängt.\\[-1mm]
\begin{tikzpicture}[node distance=12mm]
\node[qobj] (ws) {\textbf{ws : Warteschlange}\nodepart{second}anfang\\[5mm]ende};
\node[obj,right=14mm of ws.north east,anchor=north west] (k1) {\textbf{k1 : Knoten}\nodepart{second}daten = Resi\\nachfolger};
\node[obj,right=of k1] (k2) {\textbf{k2 : Knoten}\nodepart{second}daten = Bene\\nachfolger};
\node[obj,right=of k2] (k3) {\textbf{k3 : Knoten}\nodepart{second}daten = Gustav\\nachfolger};
\begin{scope}[lsgbild]
  \draw[zeiger] ([yshift=2.4mm]ws.east) -- ([yshift=2.4mm]ws.east -| k1.west);
  \draw[zeiger] ([yshift=-3.6mm]ws.east) -| (k3.south);
  \draw[zeiger] ([yshift=-2.5mm]k1.east) -- ([yshift=-2.5mm]k2.west);
  \draw[zeiger] ([yshift=-2.5mm]k2.east) -- ([yshift=-2.5mm]k3.west);
\end{scope}
\node[nullref,red!75!black,right=1mm of k3.east,yshift=-2.5mm] {\lsgtext{null}};
\node[font=\scriptsize,red!75!black,above=0.5mm of k1] {\lsgtext{poll() nimmt vorne weg}};
\node[font=\scriptsize,red!75!black,above=0.5mm of k3] {\lsgtext{add() hängt hinten an}};
\end{tikzpicture}\\[0.5mm]
\textbf{b)} Ergänze die Datentypen und zeichne die Beziehungen mit Namen und Multiplizitäten ein.\\[1.5mm]
\begin{tikzpicture}
\node[draw,rectangle split,rectangle split parts=3,font=\scriptsize,align=left,fill=teal!6,anchor=north] (W) at (0,0)
  {\textbf{Warteschlange}\nodepart{second}- anfang: \lsg[1.3cm]{Knoten}\\- ende: \lsg[1.3cm]{Knoten}
   \nodepart{third}+ add(Patient p): boolean\\+ poll(): Patient\\+ peek(): Patient\\+ isEmpty(): boolean\\+ ausgeben()};
\node[draw,rectangle split,rectangle split parts=3,font=\scriptsize,align=left,fill=orange!8,anchor=north] (K) at (6.9,0)
  {\textbf{Knoten}\nodepart{second}- daten: \lsg[1.3cm]{Patient}\\- nachfolger: \lsg[1.3cm]{Knoten}
   \nodepart{third}+ Knoten(Patient daten)\\+ getDaten(): Patient\\+ getNachfolger(): Knoten\\+ setNachfolger(Knoten k)};
\node[draw,rectangle split,rectangle split parts=3,font=\scriptsize,align=left,fill=gray!8,anchor=north] (P) at (13.2,0)
  {\textbf{Patient}\nodepart{second}- name: String\\- beschwerden: String
   \nodepart{third}+ getName(): String\\+ getBeschwerden(): String};
\begin{scope}[font=\scriptsize,red!75!black]
  \draw[zeiger,lsgbild] ([yshift=-5mm]W.north east) -- ([yshift=-5mm]W.north east -| K.west);
  \draw[zeiger,lsgbild] ([yshift=-11mm]W.north east) -- ([yshift=-11mm]W.north east -| K.west);
  \draw[zeiger,lsgbild] ([yshift=-8mm]K.north east) -- ([yshift=-8mm]K.north east -| P.west);
  \draw[zeiger,lsgbild] ([xshift=-6mm]K.north east) -- ++(0,5mm) -| ([xshift=6mm]K.north west);
  \node[above] at ($([yshift=-5mm]W.north east)!0.45!([yshift=-5mm]W.north east -| K.west)$) {\lsgtext{anfang}};
  \node[above left] at ([yshift=-5mm]W.north east -| K.west) {\lsgtext{0..1}};
  \node[below] at ($([yshift=-11mm]W.north east)!0.45!([yshift=-11mm]W.north east -| K.west)$) {\lsgtext{ende}};
  \node[below left] at ([yshift=-11mm]W.north east -| K.west) {\lsgtext{0..1}};
  \node[above] at ($([yshift=-8mm]K.north east)!0.45!([yshift=-8mm]K.north east -| P.west)$) {\lsgtext{daten}};
  \node[above left] at ([yshift=-8mm]K.north east -| P.west) {\lsgtext{1}};
  \node[right] at ($(K.north east)+(-5mm,3.5mm)$) {\lsgtext{nachfolger\quad 0..1}};
\end{scope}
\end{tikzpicture}
\end{aufgabe}

\begin{merke}[Die Warteschlange (Queue)]
\begin{itemize}
\item Eine Warteschlange speichert \lsg[2.5cm]{beliebig viele} Objekte in einer festen \lsg[2.3cm]{Reihenfolge}.
\item Die Warteschlange kennt nur den ersten (\code{anfang}) und letzten Knoten (\code{ende}), jeder Knoten seine \code{daten} und seinen \lsg[2.3cm]{Nachfolger}. Beim letzten Knoten ist er \lsg[1.3cm]{null}.
\item Eine Beziehung zwischen Objekten \textbf{derselben} Klasse heißt \lsg[3.8cm]{rekursive Beziehung}. Im Code wird sie zu einem Referenzattribut, dessen Typ die \lsg[2.6cm]{eigene Klasse} ist: \code{private Knoten nachfolger;}
\item \code{add} hängt \lsg[1.5cm]{hinten} an, \code{poll} nimmt \lsg[1.5cm]{vorne} weg: \lsg[1.4cm]{FIFO}, First In -- First Out.
\end{itemize}
\end{merke}

\newpage
\begin{aufgabe}[Projekt Praxis, TODO 3]{3}
\code{add(Patient p)} braucht drei Schritte. Übersetze jeden Schritt im Bild in eine Zeile Code.\\[1mm]
\begin{minipage}[t]{0.47\linewidth}
\vspace{0pt}
\begin{tikzpicture}[node distance=6mm]
\node[obj] (a) {\textbf{Knoten}\nodepart{second}Resi};
\node[obj,right=of a] (b) {\textbf{Knoten}\nodepart{second}Bene};
\node[obj,right=11mm of b] (n) {\textbf{Knoten}\nodepart{second}Gustav};
\node[font=\scriptsize\ttfamily,above=5mm of a] (an) {anfang};
\node[font=\scriptsize\ttfamily,above=5mm of b] (en) {ende};
\draw[zeiger] (an) -- (a); \draw[zeiger] (a) -- (b);
\draw[zeiger,gray,densely dashed] (en) -- (b);
\draw[gray] ($(en)!0.55!(b)+(-1.2mm,-1.2mm)$) -- ++(2.4mm,2.4mm);
\draw[zeiger,teal!60!black] (b) -- node[schritt,above=0.5mm] {2} (n);
\draw[zeiger,teal!60!black] (en) -- node[schritt,above right=-0.6mm] {3} (n.north west);
\node[schritt,below=1mm of n] {1};
\node[font=\scriptsize,below=5mm of n.south,align=center] {\code{neu}};
\end{tikzpicture}\\[1mm]
\scriptsize\raggedright
\nr{1} Neuen Knoten für \code{p} erzeugen.\\
\nr{2} Bisher letzter Knoten bekommt \code{neu} als Nachfolger.\\
\nr{3} \code{ende} verweist auf den neuen Knoten.
\end{minipage}\hfill
\begin{minipage}[t]{0.51\linewidth}
\vspace{0pt}
\begin{java}[basicstyle=\ttfamily\scriptsize]
public boolean add(Patient p) {
   Knoten neu = §\lsg[2.4cm]{new Knoten(p)}§;   // §\nr{1}§
   if (isEmpty()) {         // Sonderfall
      anfang = §\lsg[1.1cm]{neu}§;
      ende = §\lsg[1.1cm]{neu}§;
   } else {
      §\lsg[3.4cm]{ende.setNachfolger(neu)}§;     // §\nr{2}§
      §\lsg[2cm]{ende = neu}§;                    // §\nr{3}§
   }
   return true;
}
\end{java}
\end{minipage}\\[1mm]
\textbf{b)} Warum braucht man den Sonderfall für die leere Schlange? Was passiert ohne ihn?\\ \lsg[0.99\linewidth]{ende ist null: ende.setNachfolger(neu) führt zu einer NullPointerException.}\\[1mm]
\textbf{c)} Übertrage die Methode ins Projekt und teste sie mit \code{new Warteschlange()}.
\end{aufgabe}

\begin{aufgabe}[Projekt Praxis, TODO 4]{4}
\code{poll()} holt den vordersten Patienten heraus und gibt ihn zurück.\\[1mm]
\begin{minipage}[t]{0.47\linewidth}
\vspace{0pt}
\begin{tikzpicture}[node distance=6mm]
\node[obj] (a) {\textbf{Knoten}\nodepart{second}Resi};
\node[obj,right=of a] (b) {\textbf{Knoten}\nodepart{second}Bene};
\node[obj,right=of b] (c) {\textbf{Knoten}\nodepart{second}Gustav};
\node[font=\scriptsize\ttfamily,above=5mm of a] (an) {anfang};
\node[font=\scriptsize\ttfamily,above=5mm of c] (en) {ende};
\node[font=\scriptsize\ttfamily,below=5mm of a] (p) {p};
\draw[zeiger,gray,densely dashed] (an) -- (a);
\draw[gray] ($(an)!0.5!(a)+(-1.2mm,-1.2mm)$) -- ++(2.4mm,2.4mm);
\draw[zeiger] (a) -- (b); \draw[zeiger] (b) -- (c); \draw[zeiger] (en) -- (c);
\draw[zeiger,teal!60!black] (p) -- node[schritt,left=0.5mm] {1} (a);
\draw[zeiger,teal!60!black] (an) -- node[schritt,above right=-0.6mm] {2} (b.north west);
\end{tikzpicture}\\[1mm]
\scriptsize\raggedright
\nr{1} Die Daten des ersten Knotens in \code{p} merken.\\
\nr{2} \code{anfang} rückt auf den Nachfolger weiter.\\
\nr{3} War das der letzte Knoten, wird auch \code{ende} leer.\\
\nr{4} \code{p} zurückgeben.
\end{minipage}\hfill
\begin{minipage}[t]{0.51\linewidth}
\vspace{0pt}
\begin{java}[basicstyle=\ttfamily\scriptsize]
public Patient poll() {
   if (isEmpty()) {
      return §\lsg[1.1cm]{null}§;          // niemand wartet
   } else {
      Patient p = §\lsg[2.8cm]{anfang.getDaten()}§;      // §\nr{1}§
      anfang = §\lsg[3.4cm]{anfang.getNachfolger()}§;    // §\nr{2}§
      if (§\lsg[2.4cm]{anfang == null}§) {               // §\nr{3}§
         ende = null;
      }
      return §\lsg[0.8cm]{p}§;                           // §\nr{4}§
   }
}
\end{java}
\end{minipage}\\[1mm]
\textbf{b)} Warum muss Schritt 1 vor Schritt 2 stehen?\\ \lsg[0.99\linewidth]{Nach Schritt 2 verweist nichts mehr auf Resis Knoten: Man käme nicht mehr an ihre Daten.}\\[1mm]
\textbf{c)} Was geschieht danach mit dem Knoten von Resi?\\ \lsg[0.99\linewidth]{Kein Verweis zeigt mehr auf ihn: Java räumt ihn automatisch weg (Garbage Collector).}
\end{aufgabe}

\begin{aufgabe}[Projekt Praxis]{5}
Stelle das Hauptprogramm auf \code{new Warteschlange()} um und vergleiche mit Variante 3.\\[1mm]
\small\renewcommand{\arraystretch}{1.15}
\begin{tabularx}{\linewidth}{|l|X|X|}\hline
 & \textbf{Variante 3 (Aufrücken)} & \textbf{Warteschlange}\\\hline
Verschiebungen bei \code{poll()}, wenn $n$ warten\zeile & \lsgoder{\lsgfarbe $n-1$: Der Aufwand wächst mit der Länge.}{} & \lsgoder{\lsgfarbe keine: immer gleich viele Anweisungen}{}\\\hline
Wie viele können warten?\zeile & \lsgoder{\lsgfarbe höchstens 5 (Größe des Arrays)}{} & \lsgoder{\lsgfarbe beliebig viele (nur der Arbeitsspeicher begrenzt)}{}\\\hline
\end{tabularx}
\end{aufgabe}

\begin{hinweis}\small
\textbf{Es geht auch anders -- besprecht mündlich, warum man es so machen kann:}
\begin{itemize}
\item Die Warteschlange merkt sich \textbf{nur \code{anfang}}, kein \code{ende}. Wie findet \code{add} dann den letzten Knoten? Was wird einfacher, was langsamer?
\item \code{poll()} gibt den \textbf{Knoten} statt des Patienten zurück. Warum ist es besser, wenn das Hauptprogramm die Klasse \code{Knoten} gar nicht kennen muss?
\item \code{Patient} bekommt \textbf{selbst} ein Attribut \code{nachfolger}, die Klasse \code{Knoten} entfällt. Was spart man, was verliert man?
\item \code{add} liefert \code{boolean}, obwohl hier immer \code{true} herauskommt. Warum ist das trotzdem sinnvoll? (Tipp: \code{LinkedList} in der Breitensuche, Klasse 11)
\end{itemize}
\end{hinweis}

\end{document}
